% answers.tex —— 2012 年·湖北卷·理科 · 参考答案与详细解析片段
% 由 tools/generate_projects.py 依来源总账自 2012/2012_数学真题_湖北理_试题.md + 2012/2012_数学真题_湖北理_解析正文.md + 2012/2012_数学真题_湖北理_答案要点.md 生成，勿手改（改源 MD 后重跑）


\begin{answer}
\gksolhead{1}{}

\textbf{题目}：

方程 $x^2+6x+13=0$ 的一个根是\mbox{（　　）}

（A）$-3+2\mathrm i$；（B）$3+2\mathrm i$；（C）$-2+3\mathrm i$；（D）$2+3\mathrm i$。

\Needspace{4\baselineskip}
\textbf{解析}

\Needspace{5\baselineskip}
对二次三项式配方：
\nopagebreak[4]
\[\begin{aligned}
x^2+6x+13&=(x+3)^2+4,\\
(x+3)^2&=-4.
\end{aligned}\]
\Needspace{5\baselineskip}
在复数范围内开方得 $x+3=\pm2\mathrm i$，即
\nopagebreak[4]
\[x=-3\pm2\mathrm i.\]
两个根互为共轭，选项中只有 $-3+2\mathrm i$ 出现。

选 A。

\gkanswer{答案：A。配方得 $(x+3)^2=-4$，故 $x=-3\pm2\mathrm i$。}

\end{answer}
\begin{answer}
\gksolhead{2}{}

\textbf{题目}：

命题“$\exists x_0\in\complement_{\mathbf R}\mathbf Q,\ x_0^3\in\mathbf Q$”的否定是\mbox{（　　）}

（A）$\exists x_0\notin\complement_{\mathbf R}\mathbf Q,\ x_0^3\in\mathbf Q$；（B）$\exists x_0\in\complement_{\mathbf R}\mathbf Q,\ x_0^3\notin\mathbf Q$；（C）$\forall x\notin\complement_{\mathbf R}\mathbf Q,\ x^3\in\mathbf Q$；（D）$\forall x\in\complement_{\mathbf R}\mathbf Q,\ x^3\notin\mathbf Q$。

\Needspace{4\baselineskip}
\textbf{解析}

存在量词命题“$\exists x_0\in A,\ P(x_0)$”的否定是“$\forall x\in A,\ \lnot P(x)$”：量词要换，范围不能动，结论要否定。

\Needspace{5\baselineskip}
本题中 $A=\complement_{\mathbf R}\mathbf Q$，$P(x)$ 为“$x^3\in\mathbf Q$”，故否定为
\nopagebreak[4]
\[\forall x\in\complement_{\mathbf R}\mathbf Q,\quad x^3\notin\mathbf Q.\]
（A）（B）仍用存在量词，（C）把范围也否定了，都不对。

选 D。

\gkanswer{答案：D。存在量词命题的否定为全称量词命题，且要否定结论：$\forall x\in\complement_{\mathbf R}\mathbf Q$，$x^3\notin\mathbf Q$。}

\end{answer}
\begin{answer}
\gksolhead{3}{}

\textbf{题目}：

已知二次函数 $y=f(x)$ 的图象如图所示，则它与 $x$ 轴所围成的图形的面积为\mbox{（　　）}
\nopagebreak[4]
\begin{center}% ============================================================
% f2012_hubei_science_q03_parabola.tex
% 2012 湖北卷（理科）第 3 题 二次函数 y=f(x) 的图象
% 图中抛物线开口向下，顶点 (0,1)，与 x 轴交于 (-1,0)、(1,0)
% 来源 PDF：普通高考/2012/2012湖北理.pdf（SHA-256 b84a88bc…07cc）
% 原页：第 1 页（A3 横排）第 1 栏，印面 737，第 3 题题干下方
% ============================================================
\begin{tikzpicture}[>=Stealth, line width=0.7pt, font=\small, x=1.4cm, y=1.4cm]
  \draw[->] (-2.3,0) -- (2.3,0) node[right] {$x$};
  \draw[->] (0,-1.55) -- (0,1.95) node[above] {$y$};
  \draw[thick, domain=-1.62:1.62, samples=120] plot (\x,{1-(\x)*(\x)});
  \node[above left] at (0,1) {$1$};
  \node at (-1.20,-0.20) {$-1$};
  \node at (1.20,-0.20) {$1$};
  \node[below left] at (0,0) {$O$};
\end{tikzpicture}
\end{center}
\nopagebreak[4]
（A）$\dfrac{2\pi}{5}$；（B）$\dfrac43$；（C）$\dfrac32$；（D）$\dfrac{\pi}{2}$。

\Needspace{4\baselineskip}
\textbf{解析}

\Needspace{5\baselineskip}
由图读出抛物线顶点为 $(0,1)$，与 $x$ 轴交于 $(-1,0)$ 与 $(1,0)$，开口向下，故
\nopagebreak[4]
\[f(x)=a(x-1)(x+1),\qquad f(0)=-a=1,\]
\Needspace{5\baselineskip}
即
\nopagebreak[4]
\[f(x)=1-x^2.\]

\Needspace{5\baselineskip}
所求面积为曲线在 $[-1,1]$ 上与 $x$ 轴围成的部分：
\nopagebreak[4]
\[\begin{aligned}
S&=\int_{-1}^{1}(1-x^2)\,\mathrm dx\\
&=\left[x-\frac{x^3}{3}\right]_{-1}^{1}\\
&=\left(1-\frac13\right)-\left(-1+\frac13\right)\\
&=\frac43.
\end{aligned}\]

选 B。

\gkanswer{答案：B。由图象得顶点 $(0,1)$、零点 $\pm1$，即 $f(x)=1-x^2$，面积为 $\displaystyle\int_{-1}^{1}(1-x^2)\,\mathrm dx=\dfrac43$。}

\end{answer}
\begin{answer}
\gksolhead{4}{}

\textbf{题目}：

已知某几何体的三视图如图所示，则该几何体的体积为\mbox{（　　）}
\nopagebreak[4]
\begin{center}% ============================================================
% f2012_hubei_science_q04_views.tex
% 2012 湖北卷（理科）第 4 题 几何体三视图
% 正视图：直角梯形，底宽 2，左高 2，右高 4；侧视图：2×4 矩形内切半径 1 的圆；
% 俯视图：半径 1 的圆。几何体为半径 1、高 4 的圆柱被过轴心高度的斜平面所截。
% 来源 PDF：普通高考/2012/2012湖北理.pdf（SHA-256 b84a88bc…07cc）
% 原页：第 1 页（A3 横排）第 1 栏，印面 737，第 4 题题干下方
% ============================================================
\begin{tikzpicture}[>=Stealth, line width=0.7pt, font=\small, x=1cm, y=1cm]
  % ---------- 正视图 ----------
  \draw (0,0) -- (2,0) -- (2,4) -- (0,2) -- cycle;
  % 底宽 2
  \draw[<->] (0,-0.45) -- (2,-0.45);
  \draw (0,-0.58) -- (0,-0.32);
  \draw (2,-0.58) -- (2,-0.32);
  \node at (1,-0.85) {$2$};
  % 左高 2
  \draw[<->] (-0.45,0) -- (-0.45,2);
  \draw (-0.58,0) -- (-0.32,0);
  \draw (-0.58,2) -- (-0.32,2);
  \node at (-0.85,1) {$2$};
  % 右高 4
  \draw[<->] (2.45,0) -- (2.45,4);
  \draw (2.32,0) -- (2.58,0);
  \draw (2.32,4) -- (2.58,4);
  \node at (2.85,2) {$4$};
  \node at (1,-1.35) {正视图};

  % ---------- 侧视图 ----------
  \begin{scope}[xshift=4.1cm]
    \draw (0,0) rectangle (2,4);
    \draw (1,3) circle (1);
    \draw[<->] (0,-0.45) -- (2,-0.45);
    \draw (0,-0.58) -- (0,-0.32);
    \draw (2,-0.58) -- (2,-0.32);
    \node at (1,-0.85) {$2$};
    \node at (1,-1.35) {侧视图};
  \end{scope}

  % ---------- 俯视图 ----------
  \begin{scope}[yshift=-3.55cm]
    \draw (1,0) circle (1);
    \node at (1,-1.30) {俯视图};
  \end{scope}
\end{tikzpicture}
\end{center}
\nopagebreak[4]
（A）$\dfrac{8\pi}{3}$；（B）$3\pi$；（C）$\dfrac{10\pi}{3}$；（D）$6\pi$。

\Needspace{4\baselineskip}
\textbf{解析}

俯视图是半径 $1$ 的圆，说明几何体在 $xy$ 平面上的投影就是该圆；正视图是左端高 $2$、右端高 $4$ 的直角梯形，说明上表面是高度随 $x$ 线性增加的一个平面；侧视图中半径 $1$ 的圆，则是上表面边界（椭圆）在 $yz$ 平面上的投影。

\Needspace{5\baselineskip}
把底面圆取为 $x^2+y^2\leqslant1$（$z=0$），正视图沿 $y$ 方向投影。正视图左端高 $2$、右端高 $4$，故截平面为
\nopagebreak[4]
\[z=3+x,\qquad -1\leqslant x\leqslant1.\]
\Needspace{5\baselineskip}
该平面与圆柱的交线（椭圆）投影到侧视方向（$yz$ 平面）时，参数式
\nopagebreak[4]
\[y=\pm\sqrt{1-x^2},\qquad z=3+x=3\pm\sqrt{1-y^2}\]
恰好描出圆 $y^2+(z-3)^2=1$，与侧视图吻合；因此几何体正是半径 $1$、被该斜平面截得的圆柱。

\Needspace{5\baselineskip}
体积用截面高度积分：
\nopagebreak[4]
\[\begin{aligned}
V&=\iint_{x^2+y^2\leqslant1}(3+x)\,\mathrm dA\\
&=3\iint_{x^2+y^2\leqslant1}\mathrm dA+\iint_{x^2+y^2\leqslant1}x\,\mathrm dA\\
&=3\pi+0=3\pi .
\end{aligned}\]
（第二项由区域关于 $y$ 轴对称、被积函数关于 $x$ 为奇函数而为零。）

选 B。

\gkanswer{答案：B。俯视图是半径 $1$ 的圆、正视图左右高为 $2$ 与 $4$，几何体是半径 $1$ 的圆柱被过轴心高度的斜平面所截，体积 $\pi\cdot\dfrac{2+4}{2}=3\pi$。}

\end{answer}
\begin{answer}
\gksolhead{5}{}

\textbf{题目}：

设 $a\in\mathbf Z$，且 $0\leqslant a<13$，若 $5^{2012}+a$ 能被 $13$ 整除，则 $a=$\mbox{（　　）}

（A）$0$；（B）$1$；（C）$11$；（D）$12$。

\Needspace{4\baselineskip}
\textbf{解析}

\Needspace{5\baselineskip}
先算 $5$ 的幂模 $13$ 的规律：
\nopagebreak[4]
\[5^2=25=2\times13-1\equiv-1\pmod{13}.\]
\Needspace{5\baselineskip}
因为 $2012=2\times1006$，所以
\nopagebreak[4]
\[5^{2012}=(5^2)^{1006}\equiv(-1)^{1006}=1\pmod{13}.\]
\Needspace{5\baselineskip}
整除条件 $5^{2012}+a\equiv0\pmod{13}$ 于是化为
\nopagebreak[4]
\[1+a\equiv0\pmod{13},\]
\Needspace{5\baselineskip}
即 $a\equiv-1\equiv12\pmod{13}$。结合 $0\leqslant a<13$ 得
\nopagebreak[4]
\[a=12.\]

选 D。

\gkanswer{答案：D。$5^2\equiv-1\pmod{13}$，故 $5^{2012}\equiv1$，从而 $a\equiv-1\equiv12\pmod{13}$，结合 $0\leqslant a<13$ 得 $a=12$。}

\end{answer}
\begin{answer}
\gksolhead{6}{}

\textbf{题目}：

设 $a$，$b$，$c$，$x$，$y$，$z$ 是正数，且 $a^2+b^2+c^2=10$，$x^2+y^2+z^2=40$，$ax+by+cz=20$，则 $\dfrac{a+b+c}{x+y+z}=$\mbox{（　　）}

（A）$\dfrac14$；（B）$\dfrac13$；（C）$\dfrac12$；（D）$\dfrac34$。

\Needspace{4\baselineskip}
\textbf{解析}

\Needspace{5\baselineskip}
对两组数用柯西不等式：
\nopagebreak[4]
\[(ax+by+cz)^2\leqslant(a^2+b^2+c^2)(x^2+y^2+z^2)=10\times40=400.\]
而题设 $ax+by+cz=20$，其平方恰为 $400$，即不等式取等号。

\Needspace{5\baselineskip}
柯西不等式取等的条件是两向量成比例：存在实数 $\lambda$ 使
\nopagebreak[4]
\[(a,b,c)=\lambda(x,y,z).\]
\Needspace{5\baselineskip}
两边与 $(x,y,z)$ 作数量积得
\nopagebreak[4]
\[\lambda=\frac{ax+by+cz}{x^2+y^2+z^2}=\frac{20}{40}=\frac12.\]
\Needspace{5\baselineskip}
于是 $a+b+c=\lambda(x+y+z)$，即
\nopagebreak[4]
\[\frac{a+b+c}{x+y+z}=\lambda=\frac12.\]

选 C。

\gkanswer{答案：C。由柯西不等式 $(ax+by+cz)^2=(a^2+b^2+c^2)(x^2+y^2+z^2)$ 取等，得 $(a,b,c)=\dfrac12(x,y,z)$，比值为 $\dfrac12$。}

\end{answer}
\begin{answer}
\gksolhead{7}{}

\textbf{题目}：

定义在 $(-\infty,0)\cup(0,+\infty)$ 上的函数 $f(x)$，如果对于任意给定的等比数列 $\{a_n\}$，$\{f(a_n)\}$ 仍是等比数列，则称 $f(x)$ 为“保等比数列函数”。现有定义在 $(-\infty,0)\cup(0,+\infty)$ 上的如下函数：

① $f(x)=x^2$；② $f(x)=2^x$；③ $f(x)=\sqrt{|x|}$；④ $f(x)=\ln|x|$。

则其中是“保等比数列函数”的 $f(x)$ 的序号为\mbox{（　　）}

（A）①②；（B）③④；（C）①③；（D）②④。

\Needspace{4\baselineskip}
\textbf{解析}

设 $\{a_n\}$ 是公比为 $q$ 的等比数列，各项不为零。逐个检验 $\dfrac{f(a_{n+1})}{f(a_n)}$ 是否为与 $n$ 无关的常数。

\Needspace{5\baselineskip}
① $f(x)=x^2$：
\nopagebreak[4]
\[\frac{f(a_{n+1})}{f(a_n)}=\frac{a_{n+1}^2}{a_n^2}=q^2,\]
是常数，$\{f(a_n)\}$ 是等比数列。

\Needspace{5\baselineskip}
② $f(x)=2^x$：
\nopagebreak[4]
\[\frac{f(a_{n+1})}{f(a_n)}=2^{a_{n+1}-a_n}=2^{a_1q^{n-1}(q-1)},\]
当 $q\neq1$ 时随 $n$ 变化，不是常数。例如取 $a_n=2^n$，则 $f(a_n)=2^{2^n}$，相邻两项之比为 $2^{2^n}$，不是定值。

\Needspace{5\baselineskip}
③ $f(x)=\sqrt{|x|}$：
\nopagebreak[4]
\[\frac{f(a_{n+1})}{f(a_n)}=\sqrt{\frac{|a_{n+1}|}{|a_n|}}=\sqrt{|q|},\]
是常数。

④ $f(x)=\ln|x|$：取 $a_n=2^n$，则 $f(a_n)=\ln2^n=n\ln2$，是等差数列而非等比数列（$n=1$ 时该项为 $0$，更不可能成等比数列）。

故①③符合要求。

选 C。

\gkanswer{答案：C。$f(a_n)=a_n^2$ 与 $f(a_n)=\sqrt{|a_n|}$ 的相邻项之比分别为 $q^2$ 与 $\sqrt{|q|}$，仍为等比数列；$2^{a_n}$ 与 $\ln|a_n|$ 不具备此性质。}

\end{answer}
\begin{answer}
\gksolhead{8}{}

\textbf{题目}：

如图，在圆心角为直角的扇形 $OAB$ 中，分别以 $OA$，$OB$ 为直径作两个半圆。在扇形 $OAB$ 内随机取一点，则此点取自阴影部分的概率是\mbox{（　　）}
\nopagebreak[4]
\begin{center}% ============================================================
% f2012_hubei_science_q08_sector.tex
% 2012 湖北卷（理科）第 8 题 直角扇形 OAB 内的两个半圆
% 扇形圆心角 90°，OA 水平、OB 竖直；以 OA、OB 为直径各作一个半圆。
% 阴影＝两半圆的公共部分（透镜形）＋扇形内两半圆之外的部分（月牙形）。
% 来源 PDF：普通高考/2012/2012湖北理.pdf（SHA-256 b84a88bc…07cc）
% 原页：第 1 页（A3 横排）第 2 栏上部，印面 737，第 8 题题干下方
% ============================================================
\begin{tikzpicture}[>=Stealth, line width=0.7pt, font=\small, scale=2.3]
  \def\r{1}
  \coordinate (O) at (0,0);
  \coordinate (A) at (\r,0);
  \coordinate (B) at (0,\r);

  % 扇形内先铺阴影，再用白色挖去两个半圆，最后补回公共部分
  \begin{scope}
    \clip (O) -- (A) arc (0:90:\r) -- cycle;
    \fill[black!15] (O) -- (A) arc (0:90:\r) -- cycle;
    \fill[white] (\r/2,0) circle (\r/2);
    \fill[white] (0,\r/2) circle (\r/2);
  \end{scope}
  \begin{scope}
    \clip (\r/2,0) circle (\r/2);
    \clip (0,\r/2) circle (\r/2);
    \fill[black!15] (-0.2,-0.2) rectangle (\r+0.2,\r+0.2);
  \end{scope}

  % 轮廓
  \draw (O) -- (A);
  \draw (O) -- (B);
  \draw (A) arc (0:90:\r);
  \draw (O) arc (180:0:\r/2);
  \draw (B) arc (90:-90:\r/2);

  \node[below left] at (O) {$O$};
  \node[below] at (A) {$A$};
  \node[left] at (B) {$B$};
\end{tikzpicture}
\end{center}
\nopagebreak[4]
（A）$1-\dfrac{2}{\pi}$；（B）$\dfrac12-\dfrac{1}{\pi}$；（C）$\dfrac{2}{\pi}$；（D）$\dfrac{1}{\pi}$。

\Needspace{4\baselineskip}
\textbf{解析}

\Needspace{5\baselineskip}
设 $OA=OB=R$。扇形面积为 $\dfrac14\pi R^2$；以 $OA$、$OB$ 为直径的半圆半径均为 $\dfrac R2$，面积各为 $\dfrac12\pi\left(\dfrac R2\right)^2=\dfrac{\pi R^2}{8}$，两个半圆面积之和为
\nopagebreak[4]
\[\frac{\pi R^2}{8}+\frac{\pi R^2}{8}=\frac{\pi R^2}{4},\]
恰等于扇形面积。

\Needspace{5\baselineskip}
记两半圆的公共部分面积为 $S_0$，则两半圆并集面积为 $\dfrac{\pi R^2}{4}-S_0$，于是扇形内两半圆之外的部分面积为
\nopagebreak[4]
\[\frac{\pi R^2}{4}-\left(\frac{\pi R^2}{4}-S_0\right)=S_0.\]
阴影由这部分与公共部分 $S_0$ 拼成，故阴影面积为 $2S_0$。

\Needspace{5\baselineskip}
$S_0$ 是两段弓形之和。两圆半径均为 $\dfrac R2$，圆心距为 $\dfrac{\sqrt2}{2}R$，公共弦所对圆心角为 $90^\circ$，故每段弓形面积为
\nopagebreak[4]
\[\frac14\pi\left(\frac R2\right)^2-\frac12\left(\frac R2\right)^2=\frac{\pi R^2}{16}-\frac{R^2}{8},\]
\Needspace{5\baselineskip}
从而
\nopagebreak[4]
\[S_0=2\left(\frac{\pi R^2}{16}-\frac{R^2}{8}\right)=\frac{\pi R^2}{8}-\frac{R^2}{4}.\]
\Needspace{5\baselineskip}
阴影面积为
\nopagebreak[4]
\[2S_0=\frac{\pi R^2}{4}-\frac{R^2}{2},\]
\Needspace{5\baselineskip}
所求概率为几何概型
\nopagebreak[4]
\[P=\frac{\dfrac{\pi R^2}{4}-\dfrac{R^2}{2}}{\dfrac{\pi R^2}{4}}=1-\frac2\pi .\]

选 A。

\gkanswer{答案：A。两半圆面积之和恰等于扇形面积，故阴影面积为两半圆公共部分的两倍，即 $2\left(\dfrac{\pi}{8}-\dfrac14\right)R^2$，概率为 $1-\dfrac2\pi$。}

\end{answer}
\begin{answer}
\gksolhead{9}{}

\textbf{题目}：

函数 $f(x)=x\cos x^2$ 在区间 $[0,4]$ 上的零点个数为\mbox{（　　）}

（A）$4$；（B）$5$；（C）$6$；（D）$7$。

\Needspace{4\baselineskip}
\textbf{解析}

\Needspace{5\baselineskip}
$f(x)=x\cos x^2$ 的零点满足
\nopagebreak[4]
\[x=0\quad\text{\text{或}}\quad\cos x^2=0.\]
\Needspace{5\baselineskip}
后者给出
\nopagebreak[4]
\[x^2=\frac\pi2+k\pi,\qquad k\in\mathbf Z,\ k\geqslant0,\]
\Needspace{5\baselineskip}
即
\nopagebreak[4]
\[x=\sqrt{\frac\pi2+k\pi}.\]
\Needspace{5\baselineskip}
要求在 $[0,4]$ 内，即
\nopagebreak[4]
\[\begin{aligned}
\frac\pi2+k\pi\leqslant16
&\Longrightarrow\ k\leqslant\frac{16-\frac\pi2}{\pi}\\
&=\frac{16}{\pi}-\frac12\approx4.59,
\end{aligned}\]
\Needspace{5\baselineskip}
故 $k=0,1,2,3,4$，共 $5$ 个正零点：
\nopagebreak[4]
\[x\approx1.253,\ 2.171,\ 2.802,\ 3.316,\ 3.760.\]
再加上 $x=0$，零点共 $6$ 个。

选 C。

\gkanswer{答案：C。零点为 $x=0$ 及 $x=\sqrt{\dfrac\pi2+k\pi}$（$k=0,1,2,3,4$），共 $6$ 个。}

\end{answer}
\begin{answer}
\gksolhead{10}{}

\textbf{题目}：

我国古代数学名著《九章算术》中“开立圆术”曰：置积尺数，以十六乘之，九而一，所得开立方除之，即立圆径。“开立圆术”相当于给出了已知球的体积 $V$，求其直径 $d$ 的一个近似公式 $d\approx\sqrt[3]{\dfrac{16}{9}V}$。人们还用过一些类似的近似公式。根据 $\pi=3.14159\cdots$ 判断，下列近似公式中最精确的一个是\mbox{（　　）}

（A）$d\approx\sqrt[3]{\dfrac{16}{9}V}$；（B）$d\approx\sqrt[3]{2V}$；（C）$d\approx\sqrt[3]{\dfrac{300}{157}V}$；（D）$d\approx\sqrt[3]{\dfrac{21}{11}V}$。

\Needspace{4\baselineskip}
\textbf{解析}

\Needspace{5\baselineskip}
设球半径为 $r$，则
\nopagebreak[4]
\[V=\frac43\pi r^3=\frac{\pi d^3}{6}\ \Longrightarrow\ d=\sqrt[3]{\frac6\pi V},\]
\Needspace{5\baselineskip}
故四个近似公式相当于把系数 $\dfrac6\pi$ 换成其他有理数，系数越接近 $\dfrac6\pi$ 越精确。

\nopagebreak[4]
\[\frac6\pi\approx1.9098593.\]
逐一比较各公式的系数：

\begin{itemize}
\item （A）$\dfrac{16}{9}\approx1.7777778$，与 $\dfrac6\pi$ 相差约 $0.1321$；
\item （B）$2$，相差约 $0.0901$；
\item （C）$\dfrac{300}{157}\approx1.9108280$，相差约 $0.0010$；
\item （D）$\dfrac{21}{11}\approx1.9090909$，相差约 $0.0008$。
\end{itemize}

（D）的系数与 $\dfrac6\pi$ 的误差最小，故（D）最精确。

选 D。

\gkanswer{答案：D。由 $V=\dfrac{\pi d^3}{6}$ 得 $d=\sqrt[3]{\dfrac6\pi V}$，而 $\dfrac6\pi\approx1.90986$，$\dfrac{21}{11}\approx1.90909$ 与它最接近。}

\end{answer}
\begin{answer}
\gksolhead{11}{}

\textbf{题目}：

设 $\triangle ABC$ 的内角 $A$，$B$，$C$ 所对的边分别为 $a$，$b$，$c$。若 $(a+b-c)(a+b+c)=ab$，则角 $C=$\gkblank{}。

\Needspace{4\baselineskip}
\textbf{解析}

\Needspace{5\baselineskip}
把条件左边按平方差展开：
\nopagebreak[4]
\[(a+b-c)(a+b+c)=(a+b)^2-c^2=a^2+2ab+b^2-c^2.\]
\Needspace{5\baselineskip}
由题设它等于 $ab$，故
\nopagebreak[4]
\[a^2+b^2-c^2+2ab=ab,\]
\Needspace{5\baselineskip}
即
\nopagebreak[4]
\[a^2+b^2-c^2=-ab.\]
\Needspace{5\baselineskip}
由余弦定理
\nopagebreak[4]
\[\cos C=\frac{a^2+b^2-c^2}{2ab}=\frac{-ab}{2ab}=-\frac12.\]
\Needspace{5\baselineskip}
因为 $0<C<\pi$，所以
\nopagebreak[4]
\[C=\frac{2\pi}{3}.\]

\gkanswer{答案：$\dfrac{2\pi}{3}$。条件化为 $a^2+b^2-c^2=-ab$，故 $\cos C=-\dfrac12$。}

\end{answer}
\begin{answer}
\gksolhead{12}{}

\textbf{题目}：

阅读如图所示的程序框图，运行相应的程序，输出的结果 $s=$\gkblank{}。
\nopagebreak[4]
\begin{center}% ============================================================
% f2012_hubei_science_q12_flowchart.tex
% 2012 湖北卷（理科）第 12 题 程序框图
% 初值 a=1、s=0、n=1；循环体 s=s+a、a=a+2；判断 n<3，是则 n=n+1 返回，
% 否则输出 s。运行结果 s=9。
% 来源 PDF：普通高考/2012/2012湖北理.pdf（SHA-256 b84a88bc…07cc）
% 原页：第 1 页（A3 横排）第 3 栏，印面 737，第 12 题题干下方
% ============================================================
\begin{tikzpicture}[>=Stealth, line width=0.7pt, font=\small]
  % 开始
  \draw[rounded corners=8pt] (-0.55,0.50) rectangle (0.55,0.10);
  \node at (0,0.30) {开始};
  \draw[->] (0,0.10) -- (0,-0.30);

  % 初值
  \draw (-1.55,-0.30) rectangle (1.55,-0.80);
  \node at (0,-0.55) {$a=1,\ s=0,\ n=1$};
  \draw[->] (0,-0.80) -- (0,-1.20);

  % s = s + a
  \draw (-0.85,-1.20) rectangle (0.85,-1.70);
  \node at (0,-1.45) {$s=s+a$};
  \draw[->] (0,-1.70) -- (0,-2.05);

  % a = a + 2
  \draw (-0.85,-2.05) rectangle (0.85,-2.55);
  \node at (0,-2.30) {$a=a+2$};
  \draw[->] (0,-2.55) -- (0,-2.85);

  % 判断 n < 3
  \draw (0,-2.85) -- (1.05,-3.55) -- (0,-4.25) -- (-1.05,-3.55) -- cycle;
  \node at (0,-3.55) {$n<3$?};

  % 是：n = n + 1，回边
  \draw[->] (0,-4.25) -- (0,-4.60);
  \node[right] at (0.08,-4.38) {是};
  \draw (-0.80,-4.60) rectangle (0.80,-5.10);
  \node at (0,-4.85) {$n=n+1$};
  \draw (-0.80,-4.85) -- (-2.15,-4.85) -- (-2.15,-1.45);
  \draw[->] (-2.15,-1.45) -- (-0.85,-1.45);

  % 否：输出 s
  \draw[->] (1.05,-3.55) -- (2.15,-3.55);
  \node[above] at (1.30,-3.50) {否};
  \draw (2.45,-3.30) -- (3.95,-3.30) -- (3.45,-3.80) -- (1.95,-3.80) -- cycle;
  \node at (2.95,-3.55) {输出 $s$};
  \draw[->] (2.95,-3.80) -- (2.95,-4.50);

  % 结束
  \draw[rounded corners=8pt] (2.35,-4.50) rectangle (3.55,-4.95);
  \node at (2.95,-4.72) {结束};
\end{tikzpicture}
\end{center}
\nopagebreak[4]
\Needspace{4\baselineskip}
\textbf{解析}

按框图逐步执行。每轮先算 $s=s+a$、再算 $a=a+2$，然后判断 $n<3$：成立则 $n=n+1$ 后返回，不成立则输出 $s$。

\begin{itemize}
\item 第 1 轮：$s=0+1=1$，$a=1+2=3$；因 $n=1<3$ 成立，置 $n=2$；
\item 第 2 轮：$s=1+3=4$，$a=3+2=5$；因 $n=2<3$ 成立，置 $n=3$；
\item 第 3 轮：$s=4+5=9$，$a=5+2=7$；因 $n=3<3$ 不成立，退出循环。
\end{itemize}

\Needspace{5\baselineskip}
输出
\nopagebreak[4]
\[s=9.\]

\gkanswer{答案：$9$。循环依次给出 $(s,a,n)=(1,3,2)$、$(4,5,3)$、$(9,7,3)$，此时 $n<3$ 不成立，输出 $9$。}

\end{answer}
\begin{answer}
\gksolhead{13}{}

\textbf{题目}：

回文数是指从左到右读与从右到左读都一样的正整数。如 $22$，$121$，$3443$，$94249$ 等。显然 $2$ 位回文数有 $9$ 个：$11$，$22$，$33$，$\cdots$，$99$；$3$ 位回文数有 $90$ 个：$101$，$111$，$121$，$\cdots$，$191$，$202$，$\cdots$，$999$。则

（1） 4 位回文数有\gkblank{}个；

（2） $2n+1$（$n\in\mathbf N^*$）位回文数有\gkblank{}个。

\Needspace{4\baselineskip}
\textbf{解析}

\textbf{（1）}

\Needspace{5\baselineskip}
4 位回文数形如 $\overline{abba}$，其中 $a\in\{1,2,\cdots,9\}$，$b\in\{0,1,\cdots,9\}$。首位不能为 $0$，其余各位由对称性完全确定，故个数为
\nopagebreak[4]
\[9\times10=90.\]

\textbf{（2）}

\Needspace{5\baselineskip}
$2n+1$ 位回文数形如
\nopagebreak[4]
\[\overline{a_1a_2\cdots a_n\,x\,a_n\cdots a_2a_1},\]
\Needspace{5\baselineskip}
整串由前 $n+1$ 位（即 $a_1,a_2,\cdots,a_n,x$）唯一确定。首位 $a_1$ 有 $9$ 种取法，其余 $n$ 位各有 $10$ 种取法，故个数为
\nopagebreak[4]
\[9\times10^{n}.\]
当 $n=1$ 时得 $90$，与题中 3 位回文数的个数一致。

\gkanswer{答案（1）：$90$。4 位回文数形如 $\overline{abba}$，首位 $9$ 种、次位 $10$ 种，共 $9\times10=90$ 个。}

\gkanswer{答案（2）：$9\times10^n$。$2n+1$ 位回文数由前 $n+1$ 位决定，首位 $9$ 种、其余 $n$ 位各 $10$ 种。}

\end{answer}
\begin{answer}
\gksolhead{14}{}

\textbf{题目}：

如图，双曲线 $\dfrac{x^2}{a^2}-\dfrac{y^2}{b^2}=1$（$a,b>0$）的两顶点为 $A_1$，$A_2$，虚轴两端点为 $B_1$，$B_2$，两焦点为 $F_1$，$F_2$。若以 $A_1A_2$ 为直径的圆内切于菱形 $F_1B_1F_2B_2$，切点分别为 $A$，$B$，$C$，$D$。则

（1） 双曲线的离心率 $e=$\gkblank{}；

（2） 菱形 $F_1B_1F_2B_2$ 的面积 $S_1$ 与矩形 $ABCD$ 的面积 $S_2$ 的比值 $\dfrac{S_1}{S_2}=$\gkblank{}。
\nopagebreak[4]
\begin{center}% ============================================================
% f2012_hubei_science_q14_hyperbola.tex
% 2012 湖北卷（理科）第 14 题 双曲线与内切圆、菱形、矩形
% 双曲线 x^2/a^2-y^2/b^2=1，顶点 A1、A2，虚轴端点 B1、B2，焦点 F1、F2；
% 以 A1A2 为直径的圆（半径 a）内切于菱形 F1B1F2B2，切点 A、B、C、D 成矩形。
% 图中取 a=1.6，e=(1+sqrt5)/2，b=a*sqrt(e^2-1)=2.035，c=ae=2.589。
% 来源 PDF：普通高考/2012/2012湖北理.pdf（SHA-256 b84a88bc…07cc）
% 原页：第 1 页（A3 横排）第 3 栏上部，印面 737，第 14 题题干下方
% ============================================================
\begin{tikzpicture}[>=Stealth, line width=0.7pt, font=\small, x=1cm, y=1cm]
  \def\a{1.6}
  \def\b{2.035}
  \def\c{2.589}
  \def\xt{0.989}
  \def\yt{1.258}

  \draw[->] (-3.55,0) -- (3.75,0) node[right] {$x$};
  \draw[->] (0,-3.35) -- (0,3.45) node[above] {$y$};

  % 双曲线两支
  \draw[thick, domain=-2.55:2.55, samples=120] plot ({1.6*sqrt(1+(\x*\x)/4.14103)},\x);
  \draw[thick, domain=-2.55:2.55, samples=120] plot ({-1.6*sqrt(1+(\x*\x)/4.14103)},\x);

  % 菱形
  \draw (-\c,0) -- (0,\b) -- (\c,0) -- (0,-\b) -- cycle;

  % 内切圆
  \draw (\a,0) arc (0:360:\a);

  % 矩形 ABCD
  \draw (-\xt,\yt) -- (\xt,\yt) -- (\xt,-\yt) -- (-\xt,-\yt) -- cycle;

  % 切点
  \fill (\xt,\yt) circle (1.2pt) node[right] {$A$};
  \fill (-\xt,\yt) circle (1.2pt) node[left] {$B$};
  \fill (-\xt,-\yt) circle (1.2pt) node[left] {$C$};
  \fill (\xt,-\yt) circle (1.2pt) node[right] {$D$};

  % 顶点与焦点
  \fill (-\a,0) circle (1.2pt) node[below left] {$A_1$};
  \fill (\a,0) circle (1.2pt) node[below right] {$A_2$};
  \fill (0,\b) circle (1.2pt) node[right] {$B_2$};
  \fill (0,-\b) circle (1.2pt) node[right] {$B_1$};
  \fill (-\c,0) circle (1.2pt) node[below left] {$F_1$};
  \fill (\c,0) circle (1.2pt) node[below right] {$F_2$};
  \node[below left] at (0,0) {$O$};
\end{tikzpicture}
\end{center}
\nopagebreak[4]
\Needspace{4\baselineskip}
\textbf{解析}

\textbf{（1）}

设 $A_1(-a,0)$，$A_2(a,0)$，$B_1(0,-b)$，$B_2(0,b)$，$F_1(-c,0)$，$F_2(c,0)$，其中 $c^2=a^2+b^2$。以 $A_1A_2$ 为直径的圆圆心为 $O$，半径为 $a$。

\Needspace{5\baselineskip}
菱形的一条边所在直线为
\nopagebreak[4]
\[F_2B_2:\quad \frac xc+\frac yb=1,\ \text{\text{即}}\ bx+cy-bc=0.\]
\Needspace{5\baselineskip}
圆与它相切，故圆心 $O$ 到该直线的距离等于 $a$：
\nopagebreak[4]
\[\frac{bc}{\sqrt{b^2+c^2}}=a.\]
\Needspace{5\baselineskip}
两边平方：
\nopagebreak[4]
\[b^2c^2=a^2(b^2+c^2).\]
\Needspace{5\baselineskip}
由 $b^2=c^2-a^2=a^2(e^2-1)$，$b^2+c^2=a^2(e^2-1)+a^2e^2=a^2(2e^2-1)$，代入得
\nopagebreak[4]
\[a^2(e^2-1)\cdot a^2e^2=a^2\cdot a^2(2e^2-1),\]
\Needspace{5\baselineskip}
约去 $a^4$：
\nopagebreak[4]
\[e^2(e^2-1)=2e^2-1\ \Longrightarrow\ e^4-3e^2+1=0.\]
\Needspace{5\baselineskip}
解得
\nopagebreak[4]
\[e^2=\frac{3\pm\sqrt5}{2}.\]
\Needspace{5\baselineskip}
因 $e>1$ 需 $e^2>1$，而 $\dfrac{3-\sqrt5}{2}<1$，故取 $e^2=\dfrac{3+\sqrt5}{2}$，于是
\nopagebreak[4]
\[e=\sqrt{\frac{3+\sqrt5}{2}}=\frac{1+\sqrt5}{2}.\]

\textbf{（2）}

\Needspace{5\baselineskip}
菱形两条对角线长分别为 $2c$ 与 $2b$，故
\nopagebreak[4]
\[S_1=\frac12\cdot2c\cdot2b=2bc.\]

\Needspace{5\baselineskip}
设切点在第一象限为 $A(x_0,y_0)$，它是 $O$ 到直线 $bx+cy=bc$ 的垂足，由点到直线的垂足公式
\nopagebreak[4]
\[(x_0,y_0)=\frac{bc}{b^2+c^2}(b,c),\]
\Needspace{5\baselineskip}
即
\nopagebreak[4]
\[x_0=\frac{b^2c}{b^2+c^2},\qquad y_0=\frac{bc^2}{b^2+c^2}.\]
\Needspace{5\baselineskip}
由对称性，矩形 $ABCD$ 的边长为 $2x_0$ 与 $2y_0$，故
\nopagebreak[4]
\[S_2=4x_0y_0=\frac{4b^3c^3}{(b^2+c^2)^2}.\]
\Needspace{5\baselineskip}
于是
\nopagebreak[4]
\[\begin{aligned}
\frac{S_1}{S_2}&=\frac{2bc(b^2+c^2)^2}{4b^3c^3}\\
&=\frac{(b^2+c^2)^2}{2b^2c^2}\\
&=\frac{b^2+c^2}{2a^2}\qquad(\text{\text{用} }b^2c^2=a^2(b^2+c^2))\\
&=\frac{a^2(2e^2-1)}{2a^2}=\frac{2e^2-1}{2}.
\end{aligned}\]
\Needspace{5\baselineskip}
代入 $e^2=\dfrac{3+\sqrt5}{2}$ 得
\nopagebreak[4]
\[\frac{S_1}{S_2}=\frac{2+\sqrt5}{2}.\]

\gkanswer{答案（1）：$\dfrac{1+\sqrt5}{2}$。由圆心到菱形边的距离等于 $a$ 得 $b^2c^2=a^2(b^2+c^2)$，化为 $e^4-3e^2+1=0$。}

\gkanswer{答案（2）：$\dfrac{2+\sqrt5}{2}$。$S_1=2bc$，$S_2=\dfrac{4b^3c^3}{(b^2+c^2)^2}$，比值化简为 $\dfrac{2e^2-1}{2}$。}

\end{answer}
\begin{answer}
\gksolhead{15}{}

\textbf{题目}：

如图，点 $D$ 在 $\odot O$ 的弦 $AB$ 上移动，$AB=4$，连接 $OD$，过点 $D$ 作 $OD$ 的垂线交 $\odot O$ 于点 $C$，则 $CD$ 的最大值为\gkblank{}。
\nopagebreak[4]
\begin{center}% ============================================================
% f2012_hubei_science_q15_circle.tex
% 2012 湖北卷（理科）第 15 题 圆内定长弦与垂线
% 圆 O 半径 2.4，弦 AB 长 4（A、B 在圆上，弦垂直于 x 轴方向摆放）；
% D 在弦 AB 上，连 OD，过 D 作 OD 的垂线交圆于 C。
% 来源 PDF：普通高考/2012/2012湖北理.pdf（SHA-256 b84a88bc…07cc）
% 原页：第 1 页（A3 横排）第 3 栏中部，印面 737，第 15 题题干下方
% ============================================================
\begin{tikzpicture}[>=Stealth, line width=0.7pt, font=\small, x=1cm, y=1cm]
  \coordinate (O) at (0,0);
  \coordinate (A) at (-1.327,-2);
  \coordinate (B) at (-1.327,2);
  \coordinate (D) at (-1.327,1.6);
  \coordinate (C) at (-0.403,2.366);

  \draw (0,0) circle (2.4);
  \draw (A) -- (B);
  \draw (O) -- (D);
  \draw (D) -- (C);

  % 直角标记（D 处，OD 与 DC 之间）
  \draw (-1.212,1.461) -- (-1.073,1.576) -- (-1.188,1.715);

  \fill (O) circle (1.2pt) node[below] {$O$};
  \fill (A) circle (1.2pt) node[below left] {$A$};
  \fill (B) circle (1.2pt) node[above left] {$B$};
  \fill (D) circle (1.2pt) node[left] {$D$};
  \fill (C) circle (1.2pt);
  \node at (-0.06,2.62) {$C$};
\end{tikzpicture}
\end{center}
\nopagebreak[4]
\Needspace{4\baselineskip}
\textbf{解析}

\Needspace{5\baselineskip}
设圆半径为 $R$，圆心 $O$ 到弦 $AB$ 的距离为 $d_0$。由垂径定理
\nopagebreak[4]
\[d_0^2=R^2-\left(\frac{AB}{2}\right)^2=R^2-4.\]

\Needspace{5\baselineskip}
设 $D$ 到垂足的距离为 $t$（$0\leqslant t\leqslant2$），则
\nopagebreak[4]
\[OD^2=d_0^2+t^2.\]
\Needspace{5\baselineskip}
因为 $CD\perp OD$，$C$ 在圆上，故 $\triangle ODC$ 为直角三角形，$OC=R$，于是
\nopagebreak[4]
\[\begin{aligned}
CD^2&=OC^2-OD^2\\
&=R^2-\left(d_0^2+t^2\right)\\
&=R^2-\left(R^2-4\right)-t^2\\
&=4-t^2\leqslant4.
\end{aligned}\]
\Needspace{5\baselineskip}
所以 $CD\leqslant2$，当且仅当 $t=0$，即 $D$ 为 $AB$ 的中点时取到最大值
\nopagebreak[4]
\[CD_{\max}=2.\]

\gkanswer{答案：$2$。设圆心到弦的距离为 $d_0$，$D$ 到垂足的距离为 $t$，则 $CD^2=4-t^2\leqslant4$，当 $D$ 为 $AB$ 中点时取等。}

\end{answer}
\begin{answer}
\gksolhead{16}{}

\textbf{题目}：

在直角坐标系 $xOy$ 中，以原点 $O$ 为极点，$x$ 轴的正半轴为极轴建立极坐标系。已知射线 $\theta=\dfrac{\pi}{4}$ 与曲线 $\begin{cases}x=t+1,\\y=(t-1)^2\end{cases}$（$t$ 为参数）相交于 $A$，$B$ 两点，则线段 $AB$ 的中点的直角坐标为\gkblank{}。

\Needspace{4\baselineskip}
\textbf{解析}

\Needspace{5\baselineskip}
射线 $\theta=\dfrac\pi4$ 的直角坐标方程为
\nopagebreak[4]
\[y=x,\qquad x\geqslant0.\]
\Needspace{5\baselineskip}
把曲线的参数方程代入：
\nopagebreak[4]
\[(t-1)^2=t+1,\]
\Needspace{5\baselineskip}
整理得
\nopagebreak[4]
\[t^2-3t=0\ \Longrightarrow\ t=0\ \text{\text{或}}\ t=3.\]
$t=0$ 对应 $A(1,1)$，$t=3$ 对应 $B(4,4)$，两点都满足 $x\geqslant0$。

\Needspace{5\baselineskip}
线段 $AB$ 的中点坐标为
\nopagebreak[4]
\[\left(\frac{1+4}{2},\frac{1+4}{2}\right)=\left(\frac52,\frac52\right).\]

\gkanswer{答案：$\left(\dfrac52,\dfrac52\right)$。射线为 $y=x$（$x\geqslant0$），代入参数方程得 $t^2-3t=0$，交点 $(1,1)$、$(4,4)$，中点即所求。}

\end{answer}
\begin{answer}
\gksolhead{17}{}

\textbf{题目}：

已知向量 $\boldsymbol a=(\cos\omega x-\sin\omega x,\ \sin\omega x)$，$\boldsymbol b=(-\cos\omega x-\sin\omega x,\ 2\sqrt3\cos\omega x)$，设函数 $f(x)=\boldsymbol a\cdot\boldsymbol b+\lambda$（$x\in\mathbf R$）的图象关于直线 $x=\pi$ 对称，其中 $\omega$，$\lambda$ 为常数，且 $\omega\in\left(\dfrac12,1\right)$。

（1） 求函数 $f(x)$ 的最小正周期；

（2） 若 $y=f(x)$ 的图象经过点 $\left(\dfrac{\pi}{4},0\right)$，求函数 $f(x)$ 在区间 $\left[0,\dfrac{3\pi}{5}\right]$ 上的取值范围。

\Needspace{4\baselineskip}
\textbf{解析}

\Needspace{5\baselineskip}
先算数量积。由
\nopagebreak[4]
\[\boldsymbol a=(\cos\omega x-\sin\omega x,\ \sin\omega x),\qquad
\boldsymbol b=(-\cos\omega x-\sin\omega x,\ 2\sqrt3\cos\omega x),\]
\Needspace{5\baselineskip}
得
\nopagebreak[4]
\[\begin{aligned}
\boldsymbol a\cdot\boldsymbol b
&=(\cos\omega x-\sin\omega x)(-\cos\omega x-\sin\omega x)+2\sqrt3\sin\omega x\cos\omega x\\
&=-\left(\cos^2\omega x-\sin^2\omega x\right)+\sqrt3\sin2\omega x\\
&=-\cos2\omega x+\sqrt3\sin2\omega x\\
&=2\sin\left(2\omega x-\frac\pi6\right).
\end{aligned}\]
\Needspace{5\baselineskip}
故
\nopagebreak[4]
\[f(x)=2\sin\left(2\omega x-\frac\pi6\right)+\lambda.\]

\textbf{（1）}

\Needspace{5\baselineskip}
$f$ 的对称轴满足
\nopagebreak[4]
\[2\omega x-\frac\pi6=\frac\pi2+k\pi,\qquad k\in\mathbf Z,\]
\Needspace{5\baselineskip}
即
\nopagebreak[4]
\[x=\frac{\frac{2\pi}{3}+k\pi}{2\omega}.\]
\Needspace{5\baselineskip}
$x=\pi$ 是对称轴，故
\nopagebreak[4]
\[\pi=\frac{\frac{2\pi}{3}+k\pi}{2\omega}\ \Longrightarrow\ \omega=\frac13+\frac k2.\]
\Needspace{5\baselineskip}
由 $\dfrac12<\omega<1$ 得 $\dfrac16<k<\dfrac43$，整数 $k=1$，于是
\nopagebreak[4]
\[\omega=\frac56.\]
\Needspace{5\baselineskip}
最小正周期为
\nopagebreak[4]
\[T=\frac{2\pi}{2\omega}=\frac{\pi}{\omega}=\frac{6\pi}{5}.\]

\textbf{（2）}

\Needspace{5\baselineskip}
图象过点 $\left(\dfrac\pi4,0\right)$，即 $f\left(\dfrac\pi4\right)=0$。此时
\nopagebreak[4]
\[2\cdot\frac56\cdot\frac\pi4-\frac\pi6=\frac{5\pi}{12}-\frac{\pi}{6}=\frac\pi4,\]
\Needspace{5\baselineskip}
故
\nopagebreak[4]
\[2\sin\frac\pi4+\lambda=\sqrt2+\lambda=0\ \Longrightarrow\ \lambda=-\sqrt2,\]
\Needspace{5\baselineskip}
从而
\nopagebreak[4]
\[f(x)=2\sin\left(\frac53x-\frac\pi6\right)-\sqrt2.\]

\Needspace{5\baselineskip}
当 $x\in\left[0,\dfrac{3\pi}{5}\right]$ 时，令 $u=\dfrac53x-\dfrac\pi6$，则
\nopagebreak[4]
\[u\in\left[-\frac\pi6,\ \frac{5\pi}{6}\right].\]
\Needspace{5\baselineskip}
在此区间上 $\sin u$ 先增后减：在 $u=-\dfrac\pi6$ 处取最小值 $-\dfrac12$，在 $u=\dfrac\pi2$ 处取最大值 $1$，端点 $u=\dfrac{5\pi}6$ 处为 $\dfrac12$。故
\nopagebreak[4]
\[\sin u\in\left[-\frac12,\ 1\right],\]
\nopagebreak[4]
\[f(x)=2\sin u-\sqrt2\in\left[-1-\sqrt2,\ 2-\sqrt2\right].\]

\gkanswer{答案（1）：$\dfrac{6\pi}{5}$。化简得 $f(x)=2\sin\left(2\omega x-\dfrac\pi6\right)+\lambda$，由对称轴 $x=\pi$ 与 $\omega\in\left(\dfrac12,1\right)$ 得 $\omega=\dfrac56$。}

\gkanswer{答案（2）：$\left[-1-\sqrt2,\,2-\sqrt2\right]$。由过点 $\left(\dfrac\pi4,0\right)$ 得 $\lambda=-\sqrt2$，再求 $2\sin\left(\dfrac53x-\dfrac\pi6\right)-\sqrt2$ 在 $\left[0,\dfrac{3\pi}5\right]$ 上的值域。}

\end{answer}
\begin{answer}
\gksolhead{18}{}

\textbf{题目}：

已知等差数列 $\{a_n\}$ 前三项的和为 $-3$，前三项的积为 $8$。

（1） 求等差数列 $\{a_n\}$ 的通项公式；

（2） 若 $a_2$，$a_3$，$a_1$ 成等比数列，求数列 $\{|a_n|\}$ 的前 $n$ 项和。

\Needspace{4\baselineskip}
\textbf{解析}

\textbf{（1）}

\Needspace{5\baselineskip}
设公差为 $d$。前三项为 $a_1$，$a_1+d$，$a_1+2d$，其和
\nopagebreak[4]
\[3a_1+3d=3a_2=-3\ \Longrightarrow\ a_2=-1.\]
\Needspace{5\baselineskip}
前三项之积
\nopagebreak[4]
\[a_1a_2a_3=8\ \Longrightarrow\ a_1a_3=\frac{8}{-1}=-8.\]
\Needspace{5\baselineskip}
又 $a_1+a_3=2a_2=-2$，故 $a_1$，$a_3$ 是方程
\nopagebreak[4]
\[t^2+2t-8=0\]
的两根，解得 $t=2$ 或 $t=-4$。

\Needspace{5\baselineskip}
若 $(a_1,a_3)=(2,-4)$，则 $d=-3$，
\nopagebreak[4]
\[a_n=2-3(n-1)=5-3n;\]
\Needspace{5\baselineskip}
若 $(a_1,a_3)=(-4,2)$，则 $d=3$，
\nopagebreak[4]
\[a_n=-4+3(n-1)=3n-7.\]

\textbf{（2）}

\Needspace{5\baselineskip}
若 $a_2$，$a_3$，$a_1$ 成等比数列，则
\nopagebreak[4]
\[a_3^2=a_1a_2.\]

对 $a_n=5-3n$：$a_1=2$，$a_2=-1$，$a_3=-4$，$a_3^2=16$，而 $a_1a_2=-2$，不成立。

对 $a_n=3n-7$：$a_1=-4$，$a_2=-1$，$a_3=2$，$a_3^2=4=a_1a_2$，成立。

\Needspace{5\baselineskip}
故 $a_n=3n-7$，其各项为
\nopagebreak[4]
\[-4,\ -1,\ 2,\ 5,\ 8,\ \cdots\]
\Needspace{5\baselineskip}
负数项只有前两项，故
\nopagebreak[4]
\[S_1=|-4|=4,\qquad S_2=4+1=5.\]
\Needspace{5\baselineskip}
当 $n\geqslant2$ 时，
\nopagebreak[4]
\[\begin{aligned}
S_n&=5+\sum_{k=3}^{n}(3k-7)\\
&=5+3\left(\frac{n(n+1)}{2}-3\right)-7(n-2)\\
&=5+\frac{3n^2+3n-18}{2}-7n+14\\
&=\frac{3n^2-11n+20}{2}.
\end{aligned}\]
检验：$n=2$ 时 $\dfrac{12-22+20}{2}=5$，$n=3$ 时 $\dfrac{27-33+20}{2}=7$，均与直接求和一致。

\gkanswer{答案（1）：$a_n=5-3n$ 或 $a_n=3n-7$。由 $3a_2=-3$ 得 $a_2=-1$，$a_1a_3=-8$，$a_1+a_3=-2$，公差为 $\pm3$。}

\gkanswer{答案（2）：$S_1=4$，$n\geqslant2$ 时 $S_n=\dfrac{3n^2-11n+20}{2}$。等比条件筛选出 $a_n=3n-7$，其前三项为 $-4,-1,2$。}

\end{answer}
\begin{answer}
\gksolhead{19}{}

\textbf{题目}：

如图 1，$\angle ACB=45^\circ$，$BC=3$，过动点 $A$ 作 $AD\perp BC$，垂足 $D$ 在线段 $BC$ 上且异于点 $B$，连接 $AB$，沿 $AD$ 将 $\triangle ABD$ 折起，使 $\angle BDC=90^\circ$（如图 2 所示）。

（1） 当 $BD$ 的长为多少时，三棱锥 $A-BCD$ 的体积最大；

（2） 当三棱锥 $A-BCD$ 的体积最大时，设点 $E$，$M$ 分别为棱 $BC$，$AC$ 的中点，试在棱 $CD$ 上确定一点 $N$，使得 $EN\perp BM$，并求 $EN$ 与平面 $BMN$ 所成角的大小。
\nopagebreak[4]
\begin{center}% ============================================================
% f2012_hubei_science_q19_fig1.tex
% 2012 湖北卷（理科）第 19 题 图 1（折起前的平面图）
% B、D、C 共线，AD 垂直 BC，∠ACB=45°，BC=3（图中 BD=1，DC=AD=2）。
% 来源 PDF：普通高考/2012/2012湖北理.pdf（SHA-256 b84a88bc…07cc）
% 原页：第 2 页（A3 横排）第 1 栏下部，印面 738，第 19 题题干下方
% ============================================================
\begin{tikzpicture}[>=Stealth, line width=0.7pt, font=\small, x=1cm, y=1cm]
  \coordinate (B) at (0,0);
  \coordinate (D) at (1,0);
  \coordinate (C) at (3,0);
  \coordinate (A) at (1,2);

  \draw (A) -- (B) -- (C) -- cycle;
  \draw (A) -- (D);

  \fill (A) circle (1.2pt) node[above] {$A$};
  \fill (B) circle (1.2pt) node[below left] {$B$};
  \fill (C) circle (1.2pt) node[below right] {$C$};
  \fill (D) circle (1.2pt) node[below] {$D$};
\end{tikzpicture}
\end{center}
\nopagebreak[4]
\nopagebreak[4]
\begin{center}% ============================================================
% f2012_hubei_science_q19_fig2.tex
% 2012 湖北卷（理科）第 19 题 图 2（沿 AD 折起后的三棱锥 A-BCD）
% AD ⊥ 底面 BCD，∠BDC=90°；M 为 AC 中点，E 为 BC 中点；N 待定，原页未画出。
% 采用固定斜投影：D=(1.00,1.05)、B=(0.10,0.90)、C=(2.60,0.95)、A=(1.00,3.75)，
% 其中 E、M 由相应顶点直接取中点得到。
% 来源 PDF：普通高考/2012/2012湖北理.pdf（SHA-256 b84a88bc…07cc）
% 原页：第 2 页（A3 横排）第 1 栏下部，印面 738，第 19 题题干下方
% ============================================================
\begin{tikzpicture}[>=Stealth, line width=0.7pt, font=\small, x=1cm, y=1cm]
  \coordinate (A) at (1.00,3.75);
  \coordinate (B) at (0.10,0.90);
  \coordinate (C) at (2.60,0.95);
  \coordinate (D) at (1.00,1.05);
  \coordinate (E) at (1.35,0.925);
  \coordinate (M) at (1.80,2.35);

  % 可见棱
  \draw (A) -- (B);
  \draw (A) -- (C);
  \draw (B) -- (C);
  \draw (B) -- (M);

  % 被遮棱
  \draw[dashed] (A) -- (D);
  \draw[dashed] (D) -- (B);
  \draw[dashed] (D) -- (C);

  % D 处直角标记（DA 与 DC 之间）
  \draw (1.00,1.25) -- (1.20,1.2375) -- (1.20,1.0375);

  \fill (A) circle (1.2pt) node[above] {$A$};
  \fill (B) circle (1.2pt) node[below left] {$B$};
  \fill (C) circle (1.2pt) node[right] {$C$};
  \fill (D) circle (1.2pt) node[below left] {$D$};
  \fill (E) circle (1.2pt) node[below right] {$E$};
  \fill (M) circle (1.2pt) node[right] {$M$};
\end{tikzpicture}
\end{center}
\nopagebreak[4]
\Needspace{4\baselineskip}
\textbf{解析}

\textbf{（1）}

\Needspace{5\baselineskip}
在折起前的平面图（图 1）中，$AD\perp BC$，$\angle ACD=45^\circ$，故 $\triangle ADC$ 为等腰直角三角形，
\nopagebreak[4]
\[AD=DC=BC-BD=3-x,\]
其中 $x=BD$，$0<x<3$。

\Needspace{5\baselineskip}
折起后 $AD\perp BD$，$AD\perp DC$，且 $BD$ 与 $DC$ 相交于 $D$，故
\nopagebreak[4]
\[AD\perp\text{\text{平面} }BCD.\]
\Needspace{5\baselineskip}
又 $\angle BDC=90^\circ$，$\triangle BCD$ 的面积为 $\dfrac12x(3-x)$，于是
\nopagebreak[4]
\[\begin{aligned}
V&=\frac13\cdot AD\cdot S_{\triangle BCD}\\
&=\frac13(3-x)\cdot\frac12x(3-x)\\
&=\frac{x(3-x)^2}{6}.
\end{aligned}\]
\Needspace{5\baselineskip}
求导：
\nopagebreak[4]
\[V'(x)=\frac{(3-x)^2-2x(3-x)}{6}=\frac{(3-x)(3-3x)}{6}.\]
\Needspace{5\baselineskip}
在 $(0,3)$ 内 $V'=0$ 只有 $x=1$：$x<1$ 时 $V'>0$，$x>1$ 时 $V'<0$，故 $x=1$ 为最大值点，此时
\nopagebreak[4]
\[BD=1,\qquad DC=AD=2,\qquad V_{\max}=\frac{1\times4}{6}=\frac23.\]

\textbf{（2）}

\Needspace{5\baselineskip}
取 $D$ 为原点，$\overrightarrow{DB}$、$\overrightarrow{DC}$、$\overrightarrow{DA}$ 分别为 $x$、$y$、$z$ 轴正方向。此时
\nopagebreak[4]
\[D(0,0,0),\quad B(1,0,0),\quad C(0,2,0),\quad A(0,0,2),\]
\nopagebreak[4]
\[E\left(\frac12,1,0\right),\qquad M(0,1,1).\]
\Needspace{5\baselineskip}
设 $N(0,t,0)$，$0\leqslant t\leqslant2$，则
\nopagebreak[4]
\[\overrightarrow{EN}=\left(-\frac12,\ t-1,\ 0\right),\qquad \overrightarrow{BM}=(-1,1,1).\]
\Needspace{5\baselineskip}
由 $EN\perp BM$ 得
\nopagebreak[4]
\[\overrightarrow{EN}\cdot\overrightarrow{BM}=\frac12+(t-1)=t-\frac12=0,\]
\Needspace{5\baselineskip}
故
\nopagebreak[4]
\[t=\frac12,\]
即 $N$ 为棱 $CD$ 上满足 $DN=\dfrac12$（$DN=\dfrac14DC$）的点。

\Needspace{5\baselineskip}
此时 $\overrightarrow{EN}=\left(-\dfrac12,-\dfrac12,0\right)$。取 $\overrightarrow{BM}=(-1,1,1)$，$\overrightarrow{BN}=\left(-1,\dfrac12,0\right)$，设平面 $BMN$ 的法向量为 $\boldsymbol n=(n_x,n_y,n_z)$，则
\nopagebreak[4]
\[\begin{aligned}
n_x&=1\cdot0-1\cdot\tfrac12=-\tfrac12,\\
n_y&=1\cdot(-1)-(-1)\cdot0=-1,\\
n_z&=(-1)\cdot\tfrac12-1\cdot(-1)=\tfrac12,
\end{aligned}\]
即 $\boldsymbol n=\left(-\tfrac12,\ -1,\ \tfrac12\right)$。
\Needspace{5\baselineskip}
设 $EN$ 与平面 $BMN$ 所成角为 $\theta$，则
\nopagebreak[4]
\[\begin{aligned}
\sin\theta&=\frac{\left|\overrightarrow{EN}\cdot\boldsymbol n\right|}{\left|\overrightarrow{EN}\right|\left|\boldsymbol n\right|}\\
&=\frac{\left|\frac14+\frac12\right|}{\frac{\sqrt2}{2}\cdot\sqrt{\frac32}}
=\frac{\frac34}{\frac{\sqrt3}{2}}=\frac{\sqrt3}{2}.
\end{aligned}\]
故 $\theta=60^\circ$。

\gkanswer{答案（1）：$BD=1$。此时 $DC=AD=2$，$V=\dfrac16x(3-x)^2$ 在 $x=1$ 处取最大值 $\dfrac23$。}

\gkanswer{答案（2）：$N$ 为棱 $CD$ 上满足 $DN=\dfrac12$（即 $DN=\dfrac14DC$）的点；$EN$ 与平面 $BMN$ 所成角为 $60^\circ$。以 $D$ 为原点建立坐标系，由 $\overrightarrow{EN}\cdot\overrightarrow{BM}=0$ 定出 $N$，再求线面角。}

\end{answer}
\begin{answer}
\gksolhead{20}{}

\textbf{题目}：

根据以往的经验，某工程施工期间的降水量 $X$（单位：mm）对工期的影响如下表：

\begin{center}
\begin{tabular}{lllll}
\toprule
降水量 $X$ & $X<300$ & $300\leqslant X<700$ & $700\leqslant X<900$ & $X\geqslant900$ \\
\midrule
工期延误天数 $Y$ & $0$ & $2$ & $6$ & $10$ \\
\bottomrule
\end{tabular}
\end{center}

历年气象资料表明，该工程施工期间降水量 $X$ 小于 $300$，$700$，$900$ 的概率分别为 $0.3$，$0.7$，$0.9$。求：

（1） 工期延误天数 $Y$ 的均值与方差；

（2） 在降水量 $X$ 至少是 $300$ 的条件下，工期延误不超过 $6$ 天的概率。

\Needspace{4\baselineskip}
\textbf{解析}

\textbf{（1）}

\Needspace{5\baselineskip}
由 $P(X<300)=0.3$，$P(X<700)=0.7$，$P(X<900)=0.9$ 得 $Y$ 的分布：
\nopagebreak[4]
\[\begin{aligned}
P(Y=0)&=P(X<300)=0.3,\\
P(Y=2)&=P(300\leqslant X<700)=0.7-0.3=0.4,\\
P(Y=6)&=P(700\leqslant X<900)=0.9-0.7=0.2,\\
P(Y=10)&=P(X\geqslant900)=1-0.9=0.1.
\end{aligned}\]
\Needspace{5\baselineskip}
于是
\nopagebreak[4]
\[\begin{aligned}
E(Y)&=0\times0.3+2\times0.4+6\times0.2+10\times0.1\\
&=0+0.8+1.2+1=3,\\
E(Y^2)&=0^2\times0.3+2^2\times0.4+6^2\times0.2+10^2\times0.1\\
&=0+1.6+7.2+10=18.8,\\
D(Y)&=E(Y^2)-\left[E(Y)\right]^2=18.8-9=9.8.
\end{aligned}\]

\textbf{（2）}

\Needspace{5\baselineskip}
所求为条件概率
\nopagebreak[4]
\[P(Y\leqslant6\mid X\geqslant300)=\frac{P(300\leqslant X<900)}{P(X\geqslant300)}.\]
\Needspace{5\baselineskip}
其中
\nopagebreak[4]
\[\begin{aligned}
P(X\geqslant300)&=1-P(X<300)=0.7,\\
P(300\leqslant X<900)&=P(X<900)-P(X<300)=0.9-0.3=0.6.
\end{aligned}\]
\Needspace{5\baselineskip}
故
\nopagebreak[4]
\[P(Y\leqslant6\mid X\geqslant300)=\frac{0.6}{0.7}=\frac67.\]

\gkanswer{答案（1）：$E(Y)=3$，$D(Y)=9.8$。$Y$ 的分布为 $P(Y=0)=0.3$，$P(Y=2)=0.4$，$P(Y=6)=0.2$，$P(Y=10)=0.1$。}

\gkanswer{答案（2）：$\dfrac67$。条件概率为 $\dfrac{P(300\leqslant X<900)}{P(X\geqslant300)}=\dfrac{0.6}{0.7}$。}

\end{answer}
\begin{answer}
\gksolhead{21}{}

\textbf{题目}：

设 $A$ 是单位圆 $x^2+y^2=1$ 上的任意一点，$l$ 是过点 $A$ 与 $x$ 轴垂直的直线，$D$ 是直线 $l$ 与 $x$ 轴的交点，点 $M$ 在直线 $l$ 上，且满足 $|DM|=m|DA|$（$m>0$ 且 $m\neq1$）。当点 $A$ 在圆上运动时，记点 $M$ 的轨迹为曲线 $C$。

（1） 求曲线 $C$ 的方程，判断曲线 $C$ 为何种圆锥曲线，并求其焦点坐标；

（2） 过原点且斜率为 $k$ 的直线交曲线 $C$ 于 $P$，$Q$ 两点，其中 $P$ 在第一象限，它在 $y$ 轴上的射影为点 $N$，直线 $QN$ 交曲线 $C$ 于另一点 $H$。是否存在 $m$，使得对任意的 $k>0$，都有 $PQ\perp PH$？若存在，求 $m$ 的值；若不存在，请说明理由。

\Needspace{4\baselineskip}
\textbf{解析}

\textbf{（1）}

\Needspace{5\baselineskip}
设 $A(\cos\theta,\sin\theta)$，则 $D(\cos\theta,0)$，直线 $l$ 为 $x=\cos\theta$，故 $M$ 的坐标为
\nopagebreak[4]
\[M(\cos\theta,\ y_M),\qquad |DM|=|y_M|=m|\sin\theta|,\]
\Needspace{5\baselineskip}
即
\nopagebreak[4]
\[x=\cos\theta,\qquad y=\pm m\sin\theta.\]
\Needspace{5\baselineskip}
消去参数：由 $\cos\theta=x$ 及 $\dfrac{\sin\theta}{\pm1}=\dfrac{y}{m}$ 得
\nopagebreak[4]
\[x^2+\frac{y^2}{m^2}=\cos^2\theta+\sin^2\theta=1.\]
\Needspace{5\baselineskip}
故曲线
\nopagebreak[4]
\[C:\ x^2+\frac{y^2}{m^2}=1,\]
它是椭圆（$m>0$，$m\neq1$，故不是圆）。

\Needspace{5\baselineskip}
当 $m>1$ 时，长轴在 $y$ 轴上，$a^2=m^2$，$b^2=1$，半焦距 $c=\sqrt{m^2-1}$，焦点为
\nopagebreak[4]
\[\left(0,\pm\sqrt{m^2-1}\right);\]
\Needspace{5\baselineskip}
当 $0<m<1$ 时，长轴在 $x$ 轴上，$a^2=1$，$b^2=m^2$，半焦距 $c=\sqrt{1-m^2}$，焦点为
\nopagebreak[4]
\[\left(\pm\sqrt{1-m^2},\ 0\right).\]

\textbf{（2）}

\Needspace{5\baselineskip}
设 $k>0$，直线 $y=kx$ 与 $C$ 交于 $P$，$Q$。把 $y=kx$ 代入 $x^2+\dfrac{y^2}{m^2}=1$：
\nopagebreak[4]
\[x^2\left(1+\frac{k^2}{m^2}\right)=1\ \Longrightarrow\ x^2=\frac{m^2}{m^2+k^2}.\]
\Needspace{5\baselineskip}
$P$ 在第一象限，故取
\nopagebreak[4]
\[x_0=\frac{m}{\sqrt{m^2+k^2}},\qquad P(x_0,kx_0),\qquad Q(-x_0,-kx_0).\]
\Needspace{5\baselineskip}
$P$ 在 $y$ 轴上的射影为
\nopagebreak[4]
\[N(0,kx_0).\]
\Needspace{5\baselineskip}
直线 $QN$ 过 $Q(-x_0,-kx_0)$ 与 $N(0,kx_0)$，斜率为
\nopagebreak[4]
\[\frac{kx_0-(-kx_0)}{0-(-x_0)}=2k,\]
\Needspace{5\baselineskip}
方程为
\nopagebreak[4]
\[y=2kx+kx_0.\]

\Needspace{5\baselineskip}
把该直线方程代入 $C$ 求另一交点 $H$。令 $x=t$，则 $y=2kt+kx_0$，代入得
\nopagebreak[4]
\[t^2+\frac{(2kt+kx_0)^2}{m^2}=1,\]
\Needspace{5\baselineskip}
即
\nopagebreak[4]
\[\left(1+\frac{4k^2}{m^2}\right)t^2+\frac{4k^2x_0}{m^2}t+\frac{k^2x_0^2}{m^2}-1=0.\]
\Needspace{5\baselineskip}
它的两根为 $t_Q=-x_0$ 与 $t_H$。由 $x_0^2=\dfrac{m^2}{m^2+k^2}$ 得常数项
\nopagebreak[4]
\[\frac{k^2x_0^2}{m^2}-1=\frac{k^2}{m^2+k^2}-1=-\frac{m^2}{m^2+k^2},\]
\Needspace{5\baselineskip}
故两根之积
\nopagebreak[4]
\[t_Q\,t_H=\frac{-\dfrac{m^2}{m^2+k^2}}{1+\dfrac{4k^2}{m^2}}=-\frac{m^4}{(m^2+k^2)(m^2+4k^2)},\]
\Needspace{5\baselineskip}
从而
\nopagebreak[4]
\[t_H=\frac{m^3}{(m^2+4k^2)\sqrt{m^2+k^2}}=\frac{m^2}{m^2+4k^2}\,x_0.\]
\Needspace{5\baselineskip}
于是
\nopagebreak[4]
\[H\left(t_H,\ 2kt_H+kx_0\right).\]

\Needspace{5\baselineskip}
$PQ\perp PH$ 等价于
\nopagebreak[4]
\[\overrightarrow{PQ}\cdot\overrightarrow{PH}=0.\]
\Needspace{5\baselineskip}
其中 $\overrightarrow{PQ}=(-2x_0,-2kx_0)$，而
\nopagebreak[4]
\[\begin{aligned}
\overrightarrow{PH}&=\left(t_H-x_0,\ 2kt_H+kx_0-kx_0\right)\\
&=\left(t_H-x_0,\ 2kt_H\right),
\end{aligned}\]
\Needspace{5\baselineskip}
故
\nopagebreak[4]
\[\begin{aligned}
\overrightarrow{PQ}\cdot\overrightarrow{PH}
&=-2x_0\left(t_H-x_0\right)-2kx_0\cdot2kt_H\\
&=2x_0\left[x_0-t_H\left(1+2k^2\right)\right].
\end{aligned}\]
\Needspace{5\baselineskip}
因 $x_0>0$，条件化为
\nopagebreak[4]
\[x_0=t_H\left(1+2k^2\right)=\frac{m^2\left(1+2k^2\right)}{m^2+4k^2}\,x_0,\]
\Needspace{5\baselineskip}
约去 $x_0$：
\nopagebreak[4]
\[m^2+4k^2=m^2\left(1+2k^2\right)\ \Longrightarrow\ 4k^2=2m^2k^2.\]
\Needspace{5\baselineskip}
对任意 $k>0$ 都成立，故
\nopagebreak[4]
\[m^2=2,\qquad m=\sqrt2.\]
即存在 $m=\sqrt2$，使对任意 $k>0$ 都有 $PQ\perp PH$。

\gkanswer{答案（1）：$C:\ x^2+\dfrac{y^2}{m^2}=1$，是椭圆；$m>1$ 时焦点为 $\left(0,\pm\sqrt{m^2-1}\right)$，$0<m<1$ 时焦点为 $\left(\pm\sqrt{1-m^2},0\right)$。由 $|DM|=m|DA|$ 得 $y=\pm m\sin\theta$，消去参数即得。}

\gkanswer{答案（2）：存在，$m=\sqrt2$。把直线 $QN$ 与 $C$ 的第二交点 $H$ 用 $k$、$m$ 表示后，$PQ\perp PH$ 化为与 $k$ 无关的方程 $m^2=2$。}

\end{answer}
\begin{answer}
\gksolhead{22}{}

\textbf{题目}：

（1） 已知函数 $f(x)=rx-x^r+(1-r)$（$x>0$），其中 $r$ 为有理数，且 $0<r<1$，求 $f(x)$ 的最小值；

（2） 试用（1）的结果证明如下命题：设 $a_1\geqslant0$，$a_2\geqslant0$，$b_1$，$b_2$ 为正有理数。若 $b_1+b_2=1$，则 $a_1^{b_1}a_2^{b_2}\leqslant a_1b_1+a_2b_2$；

（3） 请将（2）中的命题推广到一般形式，并用数学归纳法证明你所推广的命题。

注：当 $\alpha$ 为正有理数时，有求导公式 $(x^\alpha)'=\alpha x^{\alpha-1}$。

\Needspace{4\baselineskip}
\textbf{解析}

\textbf{（1）}

\Needspace{5\baselineskip}
求导：
\nopagebreak[4]
\[f'(x)=r-rx^{r-1}=r\left(1-x^{r-1}\right).\]
\Needspace{5\baselineskip}
因 $0<r<1$，有 $r-1<0$，故 $x^{r-1}$ 在 $(0,+\infty)$ 上单调递减，且
\nopagebreak[4]
\[x^{r-1}\begin{cases}>1,&0<x<1,\\=1,&x=1,\\<1,&x>1.\end{cases}\]
\Needspace{5\baselineskip}
于是
\nopagebreak[4]
\[f'(x)\begin{cases}<0,&0<x<1,\\=0,&x=1,\\>0,&x>1.\end{cases}\]
\Needspace{5\baselineskip}
$f$ 在 $(0,1)$ 上递减、在 $(1,+\infty)$ 上递增，最小值
\nopagebreak[4]
\[f(1)=r-1+(1-r)=0.\]

\textbf{（2）}

\Needspace{5\baselineskip}
由（1），对任意 $x>0$ 有 $f(x)\geqslant f(1)=0$，即
\nopagebreak[4]
\[x^r\leqslant rx+(1-r).\]

\Needspace{5\baselineskip}
若 $a_2>0$，取 $x=\dfrac{a_1}{a_2}$、$r=b_1$，则 $1-r=1-b_1=b_2$，得
\nopagebreak[4]
\[\left(\frac{a_1}{a_2}\right)^{b_1}\leqslant b_1\cdot\frac{a_1}{a_2}+b_2.\]
\Needspace{5\baselineskip}
两边乘以正数 $a_2$：
\nopagebreak[4]
\[a_1^{b_1}a_2^{1-b_1}\leqslant a_1b_1+a_2b_2,\]
\Needspace{5\baselineskip}
即
\nopagebreak[4]
\[a_1^{b_1}a_2^{b_2}\leqslant a_1b_1+a_2b_2.\]
若 $a_2=0$，左边 $=0$，右边 $=a_1b_1\geqslant0$，不等式同样成立（$a_1=0$ 时两边都为 $0$）。

\textbf{（3）}

\Needspace{5\baselineskip}
推广形式：设 $a_i\geqslant0$，$b_i$ 为正有理数（$i=1,2,\cdots,n$），且
\nopagebreak[4]
\[\sum_{i=1}^{n}b_i=1,\]
\Needspace{5\baselineskip}
则
\nopagebreak[4]
\[\prod_{i=1}^{n}a_i^{b_i}\leqslant\sum_{i=1}^{n}b_ia_i.\]

用数学归纳法证明。

$n=2$ 时即为（2）的结论。

\Needspace{5\baselineskip}
假设 $n$ 时命题成立，考察 $n+1$ 个正有理数 $b_1,\cdots,b_{n+1}$（和为 $1$）与 $a_1,\cdots,a_{n+1}\geqslant0$。记
\nopagebreak[4]
\[B=\sum_{i=1}^{n}b_i=1-b_{n+1}>0,\]
\Needspace{5\baselineskip}
则 $\dfrac{b_1}{B},\cdots,\dfrac{b_n}{B}$ 为正有理数且和为 $1$。由归纳假设，
\nopagebreak[4]
\[\prod_{i=1}^{n}a_i^{b_i/B}\leqslant\sum_{i=1}^{n}\frac{b_i}{B}\,a_i.\]
\Needspace{5\baselineskip}
两边都是非负数，取 $B$ 次方（$t\mapsto t^B$ 在 $[0,+\infty)$ 上递增）得
\nopagebreak[4]
\[\prod_{i=1}^{n}a_i^{b_i}\leqslant\left(\sum_{i=1}^{n}\frac{b_i}{B}a_i\right)^{B}.\]
\Needspace{5\baselineskip}
再把（2）的结论用于这两个数 $\sum\limits_{i=1}^{n}\dfrac{b_i}{B}a_i$ 与 $a_{n+1}$、权 $B$ 与 $b_{n+1}$（两者和 $B+b_{n+1}=1$）：
\nopagebreak[4]
\[\begin{aligned}
&\left(\sum_{i=1}^{n}\frac{b_i}{B}a_i\right)^{B}a_{n+1}^{b_{n+1}}\\
&\qquad\leqslant B\cdot\left(\sum_{i=1}^{n}\frac{b_i}{B}a_i\right)+b_{n+1}a_{n+1}.
\end{aligned}\]
\Needspace{5\baselineskip}
两式结合即得
\nopagebreak[4]
\[\prod_{i=1}^{n+1}a_i^{b_i}\leqslant\sum_{i=1}^{n}b_ia_i+b_{n+1}a_{n+1}=\sum_{i=1}^{n+1}b_ia_i.\]
故 $n+1$ 时命题成立。由数学归纳法，推广的命题对一切 $n\geqslant2$ 成立。

当某个 $a_i=0$ 时，左边为 $0$，右边各项非负，不等式仍成立，故条件可放宽为 $a_i\geqslant0$。

\gkanswer{答案（1）：$0$。$f'(x)=r\left(1-x^{r-1}\right)$ 在 $x=1$ 变号，$f(1)=0$ 为最小值。}

\gkanswer{答案（2）：见解析版}

\gkanswer{答案（3）：推广为“设 $a_i\geqslant0$，$b_i$ 为正有理数（$i=1,2,\cdots,n$），$\sum\limits_{i=1}^n b_i=1$，则 $\prod\limits_{i=1}^n a_i^{b_i}\leqslant\sum\limits_{i=1}^n b_i a_i$”，对 $n$ 用数学归纳法证明，归纳步把前 $n$ 项的权归一化后先用归纳假设、再用 $n=2$ 的情形。}

\end{answer}